\section{Physical information in the sharp comparison}
\label{sec:actual-collar-information}

This section identifies the entropy expression in the terminal-tree
sharp comparison with physical collar mutual information and derives
the uniform pre-comparison entropy bounds from the genuine positive-prefix
input (\emph{A two-dimensional area law from a global spectral gap},
\texttt{08-scanner.tex}, lines 868--877, 900--911 and 936--941, at revision
\texttt{adc7f1241b42e322a6451854ab7e4b4c146bf78a}), without choosing any new
truncated vector, spectral set, auxiliary label, selected stage, parameter,
or terminal probability law.

\begin{theorem}[The collar partition]
    \label{thm:actual-terminal-collar-geometry}
    \lean{TNLean.PEPS.AreaLaw.Scan.CollarScan.terminalPartition_collar_geometry}
    \leanok
    Let the scanner depth be the ambient depth of the nonempty target template
    $T$, and let $8Km\leq L$.  At every completed or pre-charge terminal
    and every original band, put
    \[
        X=A\cap T,\quad U=\text{positive near},\quad
        Y=\text{middle},\quad V=\text{far}.
    \]
    These four regions are pairwise disjoint and cover the physical sites.
    The near region is $X\cup U$, the positive near-middle region is
    $U\cup Y$, and the complement of the same truncation set is contained in $V$.
\end{theorem}

\begin{proof}
    \leanok
    \uses{thm:al_scan_receiving_regions,
        thm:al_scan_nonfar_compact}
    Independently of the terminal horizon and without any entropy premise,
    the accepted receiving-region identities and positive-depth invariant
    for the status evaluator give the disjoint cover and region unions, while
    taking complements in its far-complement containment in the truncation set
    yields the outside containment in $V$.
\end{proof}

\begin{theorem}[The exact terminal identity and weighted information bound]
    \label{thm:actual-terminal-information}
    \lean{TNLean.PEPS.AreaLaw.Scan.CollarScan.terminalPartition_collar_entropy_identity,
        TNLean.PEPS.AreaLaw.Scan.CollarScan.two_mutualInformation_le_actualTerminal_average}
    \leanok
    For the same physical vector, let
    \[
        g_{\rm pre}=S(XU)+S(V)-S(Y),\qquad
        g_{\rm post}=S(U)+S(UY)-S(Y).
    \]
    Every terminal satisfies
    \[
        2S(X)+g_{\rm pre}-g_{\rm post}=I(X:V)+I(X:YV).
    \]
    If $K>0$, any given terminal mean tree, with its unchanged weights
    $w_j$, gives
    \[
        2I(X:S_0^c)\leq 2S(X)+G_{\rm pre}-G_{\rm post},\qquad
        G_\bullet=\frac1K\sum_g\sum_jw_jg_\bullet(j,g),
    \]
    where $S_0$ is the same truncation set.
\end{theorem}

\begin{proof}
    \leanok
    \uses{thm:actual-terminal-collar-geometry,
        thm:area_law_collar_leaf_identity}
    Without requiring any normalization hypothesis, the accepted pure-state
    collar identity gives the terminal equality, and its weighted
    mutual-information monotonicity theorem with $S_0^c\subseteq V$ at every
    terminal shows that both far-side informations dominate $I(X:S_0^c)$, so
    summing over the original bands and dividing by $K>0$ yields the displayed
    factor of two.
\end{proof}

\begin{theorem}[The same selected spectrum and original vector]
    \label{thm:actual-selected-collar-information}
    \lean{TNLean.PEPS.AreaLaw.Scan.CollarScan.actualCharge_two_mutualInformation_le_selectedSharpExpression}
    \leanok
    Under the preceding collar hypotheses, with $m,M,K>0$, let $\Omega_t$
    be unit, and retain its positive-mass spectral
    window $E$, of width $w$, with mass $z>0$.  Let $S'$ be the selected
    auxiliary entropy, and define $G_{\rm pre},G_{\rm post}$ from the mapped
    charge tree at the same charge stage and parameter.
    Then
    \[
        2I_{\Omega_t}(X:S_0^c)
        \leq 2S'+G_{\rm pre}-G_{\rm post}+2(w-\log z).
    \]
    For any unit original vector $\Omega$, the same right-hand side
    bounds $2I_\Omega(X:S_0^c)$ after adding
    \[
        16\sqrt{d_{\rm tr}(\Omega_t,\Omega)}\,
        \log\bigl(e\dim\mathcal H_X\bigr).
    \]
\end{theorem}

\begin{proof}
    \leanok
    \uses{thm:actual-terminal-information,
        thm:area_law_selected_schmidt_entropy_bounds,
        thm:area_law_collar_transfer}
    The accepted selected-spectrum estimate gives
    $|S'-S_{\Omega_t}(X)|\leq w-\log z$.  Its lower side costs
    exactly $2(w-\log z)$.  The accepted mutual-information continuity
    estimate has coefficient eight; multiplying by two gives sixteen.
    Only the target dimension occurs, regardless of the far-region size.
\end{proof}

\begin{theorem}[Terminal pre-entropy bounds]
    \label{thm:actual-terminal-pre-entropy}
    \lean{TNLean.PEPS.AreaLaw.Scan.CollarScan.terminalPartition_positiveEntropy_le,
        TNLean.PEPS.AreaLaw.Scan.CollarScan.terminalPartition_preEntropy_bounds,
        TNLean.PEPS.AreaLaw.Scan.CollarScan.terminalPartition_preEntropy_max_average_bounds}
    \leanok
    Assume the induced domain graph, ambient target depth, $n>0$,
    $r_0\leq D$, $1\leq D$, $4D\leq m$, $8Km\leq L$, and the
    layer-cardinality bound $n$ through depth $L$.  Suppose the same unit vector
    obeys $S(\text{positive prefix }i)\leq B_{\rm sh}$ for every
    $0\leq i\leq L$.  Put
    \[
        B=2S(X)+2B_{\rm sh}+6nD\log q.
    \]
    Every terminal through horizon $nm$ satisfies
    \[
        S(U),S(UY)\leq B_{\rm sh}+3nD\log q,
        \qquad 0\leq g_{\rm pre}\leq B.
    \]
    For $K>0$, every terminal-tree average satisfies
    $0\leq G_{\rm pre}\leq B$, and the maximum over all terminal/band
    pairs, with zero inserted, satisfies $0\leq g_{\max}\leq\max(0,B)$.
\end{theorem}

\begin{proof}
    \leanok
    \uses{thm:al_scan_physical_prefix_approximation,
        thm:actual-terminal-collar-geometry,
        thm:area_law_regional_state_identities}
    The history prefix approximation gives nested matched
    prefixes with at most $3nD$ changed sites.  The accepted disjoint-union
    entropy cost and regional dimension bound yield the two entropy estimates.
    Subadditivity on $X,U$ and $X,UY$, purity on $V$, and $S(Y)\geq0$ give
    the upper pre-entropy bound.  Subadditivity on the disjoint near and far
    regions gives its lower bound.  Nonnegative tree weights summing to one and
    division by $K>0$ preserve the bound.
\end{proof}

\begin{theorem}[The original source prefix input]
    \label{thm:source-actual-terminal-pre-entropy}
    \lean{TNLean.PEPS.AreaLaw.Scan.CollarScan.eventually_source_actualTerminal_preEntropy_bounds}
    \leanok
    Fix $q\geq1$, a positive gap parameter $g$, and a compact volume
    coefficient $C_{\rm vol}$.  At sufficiently large $n$, uniformly before
    choosing the scanner, vectors, source entropy coefficient $C_e$, or
    exponent $\eta$, the preceding bounds hold with
    \[
        B_{\rm sh}=C_e(nL^\eta+n)+n^{-247}
    \]
    whenever the original unit vector obeys
    $S_\Omega(\text{positive prefix }i)\leq C_e(nL^\eta+n)$ for all
    $0\leq i\leq L$, the compact prefix has volume at most
    $C_{\rm vol}n^2$, and the unit vector $\Omega_t$ obeys
    \[
        d_{\rm tr}(\Omega_t,\Omega)
        \leq\sqrt{2\min(n^{-1000},g/4)/g}.
    \]
\end{theorem}

\begin{proof}
    \leanok
    \uses{thm:al_positive_prefix_entropy_transfer,
        thm:actual-terminal-pre-entropy}
    Applying the accepted uniform positive-prefix entropy continuity theorem to
    the given vectors and every fixed prefix, followed by the terminal
    entropy bounds, establishes the inequalities without re-selecting any
    vector or terminal point.
\end{proof}

The original positive-prefix input remains explicit; the source safe-box
geometry must provide it.  This section does not construct that geometric
input or prove the final sublinear power estimate.
